\begin{lemma}
\label{lemma-powers}
Let $\varphi : R \to S$ be a ring map. If
\begin{enumerate}
\item for any $x \in S$ there exists $n > 0$ such that
$x^n$ is in the image of $\varphi$, and
\item $\Ker(\varphi)$ is locally nilpotent,
\end{enumerate}
then $\varphi$ induces a homeomorphism on spectra and induces residue
field extensions satisfying the equivalent conditions of
Lemma \ref{lemma-powers-field}.
\end{lemma}

\begin{proof}
Write $I=\ker\varphi$. First use only
condition (1). Every original $x\in S$
has an equation $x^n-\varphi(y)=0$
for some $n>0$ and $y\in R$. This is
a monic equation over the actual image
$R/I$, so $R/I\to S$ is an integral
injection. Lemma
\ref{lemma-integral-overring-surjective}
therefore makes $\Spec(S)\to\Spec(R/I)$
surjective. Under the original quotient
map, $\Spec(R/I)$ identifies with
the closed subset $V(I)\subseteq\Spec(R)$.

To prove injectivity, suppose primes
$\mathfrak q,\mathfrak q'$ have the
same contraction $\mathfrak p$ and
$x\in\mathfrak q\setminus\mathfrak q'$.
Choose the original equation
$x^n=\varphi(y)$ from (1). It gives
$y\in\mathfrak p$, and therefore
$x^n\in\mathfrak q'$, contradicting
primality and $x\notin\mathfrak q'$.
Interchanging the two primes proves
equality. For this same equation,
the resulting bijection sends $D(x)$
exactly onto $D(y)\cap V(I)$: a prime
contains $x$ if and only if its
contraction contains $y$. Hence it
is open for the subspace topology
on $V(I)$, and is a homeomorphism
onto that closed subset, since the
original map on spectra is continuous.
Condition (2) says precisely that
$I\subseteq\operatorname{Nil}(R)$,
so $V(I)=\Spec(R)$ and gives the
homeomorphism in the statement.

Retain corresponding original primes
$\mathfrak q$ and $\mathfrak p$.
The map $R/\mathfrak p\to S/\mathfrak q$
is injective and thus extends to the
specified residue-field embedding.
Let $x$ in the receiving residue
field be represented by $y/z$ with
$y,z\in S$, $z\notin\mathfrak q$.
Choose original $a,b\in R$ and
positive $n,m$ with
$y^n=\varphi(a)$, $z^m=\varphi(b)$.
Then $b\notin\mathfrak p$, because
its image is $z^m\notin\mathfrak q$.
The full identity in $S_{\mathfrak q}$ is
$$
(y/z)^{nm}
=\frac{(y^n)^m}{(z^m)^n}
=\frac{\varphi(a)^m}{\varphi(b)^n}
=\varphi_{\mathfrak p,\mathfrak q}(a^m/b^n).
$$
Passing to residue fields proves
that the original $x^{nm}$ belongs
to the image of $\kappa(\mathfrak p)$.
Lemma \ref{lemma-powers-field} now
applies to these exact residue fields.
This part did not require (2).

Thus (1) alone gives a homeomorphism
onto $V(\ker\varphi)$ and the same
residue-field conclusion; under (1),
(2) is equivalent to surjectivity
on all primes of $R$. No universal
homeomorphism follows from these two
conditions alone. For an exact
counterexample take the original
inclusion
$$
R=\mathbf F_2\longrightarrow
S=\mathbf F_2[\alpha]/(\alpha^2+\alpha+1).
$$
The polynomial has no root in
$\mathbf F_2$, so $S$ is a field
with elements $0,1,\alpha,\alpha^2$,
and $\alpha^3=1$. Hence every
element of $S$ has cube in $R$,
and the kernel is zero. After the
original base change $R\to S$,
the tensor ring has the comparison
$$
S\otimes_R S\simeq S[T]/(T^2+T+1),
\qquad1\otimes\alpha\longleftrightarrow T,
$$
fixing $S$ through the first factor.
The bases $1,\alpha$ and $1,T$
prove this map and its inverse.
Retain the two distinct roots
$\alpha,\alpha^2$ and the full
factorization
$T^2+T+1=(T-\alpha)(T-\alpha^2)$.
Evaluation gives
$$
S[T]/(T^2+T+1)\longrightarrow S\times S,
\qquad[H]\longmapsto(H(\alpha),H(\alpha^2)).
$$
Its inverse sends $(a,b)$ to
$$
\left[
a\frac{T-\alpha^2}{\alpha-\alpha^2}
+b\frac{T-\alpha}{\alpha^2-\alpha}
\right].
$$
The original denominators are nonzero
in $S$; evaluating verifies both
coordinates, and a polynomial zero
at both roots is divisible by the
original product of the two distinct
linear factors. These are inverse
ring maps. The base-changed spectrum
therefore has two points mapping
to the single point of $\Spec(S)$,
and cannot be a homeomorphism.
\end{proof}